\subsection{Existence of irreducible representations}

We keep the notation from the previous number.

THEOREM 3 (Raikov). --- The weak topology on Pos1(G) coincides with the topology of compact convergence.

We will first prove a few lemmas.

Lemma 5. --- Let X be a locally compact topological space and let \(\nu\) be a positive measure on X. On every bounded subset of \(\mathcal{C}_{b}(X)\) the topology induced by the weak topology \(\sigma(L^{\infty}(X), L^{1}(X))\) is coarser than the topology of compact convergence.

Let B be a bounded subset of \(\mathcal{C}_{b}(\mathrm{X})\) and let \(M \in R_{+}\) such that \(\|\varphi\|_{\infty} \leqslant M\) for all \(\varphi \in B\) . Let \(\psi \in L^{1}(\mathrm{X}, \nu)\) and \(\varepsilon > 0\) . Fix a compact subset K of X such that

\[
\int_ {\mathrm{X-K}} | \psi | d \nu <   \varepsilon .
\]

For all \(\varphi_{1}\) and \(\varphi_{2}\) in B, we then have

\[
\begin{array}{r l} & {\langle \varphi_ {1} - \varphi_ {2}, \psi \rangle = \int_ {\mathrm{X}} \psi (\varphi_ {1} - \varphi_ {2}) d \nu} \\ & {\qquad = \int_ {\mathrm{K}} \psi (\varphi_ {1} - \varphi_ {2}) d \nu + \int_ {\mathrm{X-K}} \psi (\varphi_ {1} - \varphi_ {2}) d \nu ,} \end{array}
\]

whence

\[
| \langle \varphi_ {1} - \varphi_ {2}, \psi \rangle | \leqslant \sup _ {x \in \mathrm{K}} | \varphi_ {1} (x) - \varphi_ {2} (x) | \| \psi \| _ {1} + 2 \mathrm{M} \varepsilon ,
\]

and the lemma follows.

Lemma 6. --- Let \(\psi \in \mathrm{L}^{1}(\mathrm{G})\) . Let B be a bounded subset of the Banach space \(\mathrm{L}^{\infty}(\mathrm{G})\) . The map \(\varphi \mapsto \psi * \varphi\) from B to \(\mathcal{C}_b(\mathrm{G})\) is continuous when B is equipped with the weak topology and \(\mathcal{C}_b(\mathrm{G})\) with the topology of compact convergence.

Let \(\varphi\in\mathrm{L}^{\infty}(\mathrm{G})\) and \(\eta\in\mathrm{L}^{1}(\mathrm{G})\) . The function \(\Delta^{-1}\check{\eta}\) belongs to \(\mathrm{L}^{1}(\mathrm{G})\) and \(\langle\eta,\check{\psi}\rangle=\langle\Delta^{-1}\check{\eta},\psi\rangle\) (cf. INT, VII, p. 19, § 1, n°3, formula (22)). The map \(\varphi\mapsto\check{\varphi}\) is therefore an automorphism of the space \(\mathrm{L}^{\infty}(\mathrm{G})\) endowed with the weak topology.

Let \(\varphi\in\mathrm{L}^{\infty}(\mathrm{G})\) . By INT, VIII, p. 167, § 4, no. 5, prop. 14, the function \(\psi*\varphi\) belongs to \(\mathcal{C}_{b}(\mathrm{G})\) and satisfies

\[
\begin{array}{r l} & {(\psi * \varphi) (g) = \int_ {\mathrm{G}} \psi (h) \varphi (h ^ {- 1} g) d \mu (h)} \\ & {\qquad = \int_ {\mathrm{G}} \psi (g y) \check {\varphi} (y) d \mu (y) = \langle \check {\varphi}, \pmb {\gamma} _ {\mathrm{G}} ^ {(1)} (g ^ {- 1}) \psi \rangle} \end{array}
\]

for all \(g \in G\) . It follows that the linear map \(u: \varphi \mapsto \psi * \varphi\) is a continuous map from the space \(L^{\infty}(G)\) equipped with the weak topology to \(\mathcal{C}_{b}(G)\) endowed with the topology of simple convergence.

Let \(\varphi\in\mathrm{B}\) and \((g,h)\in\mathrm{G}\times\mathrm{G}\) . By the formula above, we have

\[
| u (\varphi) (g) - u (\varphi) (h) | \leqslant \| \check {\varphi} \| _ {\infty} \| (\boldsymbol {\gamma} _ {\mathrm{G}} ^ {(1)} (g ^ {- 1}) - \boldsymbol {\gamma} _ {\mathrm{G}} ^ {(1)} (h ^ {- 1})) \psi \| _ {1}.
\]

As B is bounded and the left regular representation of G in \(L^{1}(G)\) is continuous (no. 4 of V, p. 405), this implies that \(u(B)\) is an equicontinuous subset of \(\mathcal{C}_{b}(G)\) . The assertion then follows from what precedes and from TG, X, p. 16, th. 1.

Lemma 7. --- Let \(\psi \in \mathrm{L}^1 (\mathbf{G})\) such that \(\psi \geqslant 0\) and \(\int \psi = 1\) . Let p denote the semi-norm on \(\mathcal{C}_b(\mathbf{G})\) defined by \(p(\varphi) = |\langle \varphi ,\psi \rangle |\) for all \(\varphi \in \mathcal{C}_b(\mathbf{G})\) . For every \(\varphi \in \mathrm{Pos}_1(\mathbf{G})\) , we have

\[
\| \psi * \varphi - \varphi \| _ {\infty} \leqslant \sqrt {2 p (1 - \varphi)}.
\]

Let \(\varphi\in\mathrm{Pos}_{1}(\mathrm{G})\) and \(x\in\mathrm{G}\) . According to INT, VIII, p. 167, § 4, no. 5, prop. 14, we obtain

\[
\begin{array}{r l} | \psi * \varphi (x) - \varphi (x) | = & \left| \int_ {\mathrm{G}} (\varphi (y ^ {- 1} x) - \varphi (x)) \psi (y) d \mu (y) \right| \\ & \leqslant \int_ {\mathrm{G}} | \varphi (y ^ {- 1} x) - \varphi (x) | \psi (y) d \mu (y) \\ & \leqslant \int_ {\mathrm{G}} \sqrt {2 (1 - \mathcal {R}   \varphi (y))} \psi (y) d \mu (y) \end{array}
\]

by applying the bound (1) from V, p. 446. The Cauchy-Schwarz inequality then implies

\[
\begin{array}{c} \| \psi * \varphi - \varphi \| _ {\infty} \leqslant \sqrt {2} \Bigl (\int_ {\mathrm{G}} (1 - \mathcal {R} (\varphi)) \psi   d \mu \Bigr) ^ {1 / 2} \Bigl (\int_ {\mathrm{G}} \psi   d \mu \Bigr) ^ {1 / 2} \\ \leqslant \sqrt {2} \sqrt {p (1 - \varphi)}, \end{array}
\]

whence the result.

Lemma 8. --- Let K be a topological field, and let E and F be topological vector spaces over K. Let f be a map from E into F and X a subset of E. Let \(x \in X\) . The map f \textbar{} X is continuous at x if, for every neighborhood W of 0 in F, there exist a neighborhood U of x in E and a continuous map g from X to F such that \((f - g)(\mathrm{U} \cap \mathrm{X}) \subset \mathrm{W}\) .

Let \(W_{0}\) be a neighborhood of 0 in F and let \(V_{0}\) be a symmetric neighborhood of 0 in F such that \(V_{0} + V_{0} + V_{0} \subset W_{0}\) . Let \(U_{0}\) be a neighborhood of x in E and let g be a continuous map from X into F such that \((f - g)(\mathrm{U}_{0} \cap \mathrm{X}) \subset \mathrm{V}_{0}\) .

There exists a neighborhood \(U_{1}\) of x in E such that \(g(\mathrm{U}_{1} \cap \mathrm{X}) \subset g(x) + \mathrm{V}_{0}\) . For every \(y \in U_{0} \cap U_{1} \cap X\) , we have

\[
\begin{array}{c} f (y) - f (x) = (f (y) - g (y)) + (g (y) - g (x)) + (g (x) - f (x)) \in V_{0} + V_{0} + V_{0} \subset W_{0}, \end{array}
\]

hence \(f(y) \in f(x) + W_{0}\) , whence the result.

Let us now prove Theorem 3. Denote by \(\mathrm{Pos}_{1}(\mathrm{G})_{f}\) (resp. \(\mathrm{Pos}_{1}(\mathrm{G})_{c}\) ) the set \(\mathrm{Pos}_{1}(\mathrm{G})\) endowed with the weak topology (resp. the topology of compact convergence); likewise, denote by \(\mathcal{C}_{b}(\mathrm{G})_{f}\) (resp. \(\mathcal{C}_{b}(\mathrm{G})_{c}\) ) the space \(\mathcal{C}_{b}(\mathrm{G})\) endowed with the weak topology (resp. the topology of compact convergence).

The identity map from \(\mathrm{Pos}_{1}(\mathrm{G})_{c}\) to \(\mathrm{Pos}_{1}(\mathrm{G})_{f}\) is continuous (Lemma 5). Conversely, denote by \(\iota\) the inclusion of \(\mathrm{Pos}_{1}(\mathrm{G})_{f}\) into \(\mathcal{C}_{b}(\mathrm{G})_{c}\) . Let us prove that \(\iota\) is continuous to conclude the proof.

Let \(\varphi_{1} \in \mathrm{Pos}_{1}(\mathbf{G})\) . To verify that \(\iota\) is continuous at \(\varphi_{1}\) , we apply Lemma 8. Let W be a neighborhood of 0 in \(\mathcal{C}_{b}(\mathbf{G})_{c}\) . It suffices to find a linear map \(u\) from \(\mathcal{C}_{b}(\mathbf{G})\) to \(\mathcal{C}_{b}(\mathbf{G})\) which induces, by passing to subspaces, a continuous map from \(\mathrm{Pos}_{1}(\mathbf{G})_{f}\) to \(\mathcal{C}_{b}(\mathbf{G})_{c}\) , as well as a neighborhood U of 0 in \(\mathrm{Pos}_{1}(\mathbf{G})_{f}\) so that one has \((\iota - u)(\mathrm{U} \cap \mathrm{Pos}_{1}(\mathbf{G})) \subset \mathrm{W}\) .

There exists \(\varepsilon > 0\) such that W contains the set of \(\varphi \in \mathcal{C}_b(\mathrm{G})\) which satisfy \(\| \varphi \|_{\infty} \leqslant \varepsilon\) . Since \(\varphi_1(e) = 1\) , there exists a compact neighborhood V of \(e\) in G such that

\[
\sup _ {x \in \mathrm{V}} | 1 - \varphi_ {1} (x) | \leqslant \frac {\varepsilon^ {2}}{4}.
\]

Let \(\varphi_{\mathrm{V}}\) be the characteristic function of V, and set \(\psi_{\mathrm{V}} = \mu (\mathrm{V})^{-1}\varphi_{\mathrm{V}}\) . The linear map \(u\colon \varphi \mapsto \psi_{\mathrm{V}}*\varphi\) from \(\mathrm{Pos}_1(\mathrm{G})_f\) to \(\mathcal{C}_b(\mathrm{G})_c\) is continuous (Lemma 6).

Denote by \(q_{V}\) the seminorm \(\varphi\mapsto|\langle\varphi,\psi_{V}\rangle|\) on \(\mathcal{C}_{b}(G)\) ; it is continuous for the weak topology. Since \(\psi_{V}\) is zero outside V, it follows that \(q_{\mathrm{V}}(1-\varphi_{1})\leqslant\varepsilon^{2}/4\) ; there therefore exists a neighborhood U of \(\varphi_{1}\) in \(\mathcal{C}_{b}(\mathrm{G})_{f}\) such that \(q_{\mathrm{V}}(1-\varphi)\leqslant\varepsilon^{2}/2\) for all \(\varphi\in U\) .

Let \(\varphi\in\mathrm{U}\cap\mathrm{Pos}_{1}(\mathrm{G})\) . By Lemma 7, we have

\[
\| (\iota - u) (\varphi) \| _ {\infty} \leqslant \sqrt {2 q _ {\mathrm{V}} (1 - \varphi)} \leqslant \varepsilon
\]

hence \((\iota - u)(\mathrm{U}) \subset \mathrm{W}\) . The theorem is proved.

In general, the weak topology on \(\operatorname{Pos}_{\leqslant1}(\mathbf{G})\) does not coincide with the topology of compact convergence, as shown by the example of the group R.

We recall that we denote by \(\widehat{G}\) the set of classes of irreducible unitary representations of \(G\) (V, p. 393, def. 8).

THEOREM 4 (Gelfand--Raikov). --- For every \(x \neq e\) in G, there exists an irreducible unitary representation \(\pi\) of G such that \(\pi(x)\) is not the identity endomorphism of the space \(\pi\) .

Let \(x \in G\) be such that \(\pi(x)\) is the identity for every \(\pi \in \widehat{G}\) . It follows that \(\varphi(x) = \varphi(e) = 1\) for every extreme point \(\varphi\) of \(\mathrm{Pos}_{1}(\mathrm{G})\) (prop. 11 of V, p. 444), and therefore for every \(\varphi \in \mathrm{Pos}_{1}(\mathrm{G})\) by prop. 14, since the linear form \(\varphi \mapsto \varphi(e)\) is continuous on \(\mathrm{Pos}_{1}(\mathrm{G})\) endowed with the weak topology (th. 3). But if \(x \neq e\) , there exists \(f \in L^{2}(\mathrm{G})\) of norm 1 such that \(\langle f | \boldsymbol{\gamma}_{\mathrm{G}}(x)f \rangle = 0\) (V, p. 406, lemma 3). Since the left-hand side of this equality is of the form \(\varphi(x)\) , where \(\varphi \in \mathrm{Pos}_{1}(\mathrm{G})\) , this is a contradiction.

If G is not locally compact, it is not always true that the irreducible unitary representations of G separate the points of G.

COROLLARY. --- For all elements \(g \neq h\) in G, there exists an irreducible representation \(\pi \in \widehat{G}\) and a matrix coefficient f of \(\pi\) such that \(f(g) \neq f(h)\) . In particular, the subalgebra \(\Upsilon(\mathrm{G})\) of \(\mathcal{C}_{b}(\mathrm{G})\) separates the points of G.

Let \(g \neq h\) be in G. There exists an irreducible representation \(\pi \in \widehat{G}\) such that \(\pi(g) \neq \pi(h)\) (th. 4), so there exist x and y in the space of \(\pi\) such that \(\langle x | \pi(g)y \rangle \neq \langle x | \pi(h)y \rangle\) .

